\chapter{示例、结果与验证}

\section{本章作用}

前八章已经给出了 Pebble 的问题域、方法结构、执行路径与系统骨架。本章继续回答一个更直接的问题：如果把这些设计放回真实交互中，当前到底能观察到哪些结果，又有哪些结论已经具备验证证据。本章用代表性示例、现有页面行为、接口返回结构与测试证据，区分“当前已观察事实”与“仍需后续验证的设计意图”。

\section{代表性示例}

Pebble 当前最有代表性的示例来自三条主路径：读心翻译、模拟陪练与急救呼吸。它们分别对应理解、练习与稳态恢复三个阶段，因此能够共同展示系统闭环的基本轮廓。

\subsection{示例一：读心翻译的结构化结果}

在读心翻译路径中，页面默认示例与当前组件结构都表明，系统会把一次文本输入压缩为三类核心结果。

\begin{itemize}
  \item 表面语义：对原始输入做较为中性的重述。
  \item 潜台词：对对方话语背后的压力结构或情绪意图做压缩说明。
  \item 回复建议：给出三档候选回应，并允许复制或存入练习本。
\end{itemize}

这一结果结构与前文方法章是吻合的。当前页面中，用户在得到结果后可以继续分析、复制建议，或者把选中的方案写入练习本。由此，翻译路径展示的是一种可被继续操作的结构化分析结果，包含表面语义、潜台词与三档回复建议。

\subsection{示例二：模拟陪练的回合反馈}

在模拟陪练路径中，系统返回一组随回合推进而更新的反馈对象，包括情绪得分、标签、即时反馈与注意点。当前结构中，右侧面板至少包含情绪得分、标签、即时反馈与注意点。用户每发送一轮消息，这些对象都会随着后端返回结果而更新。

这意味着陪练结果可从两个层面理解。

\begin{itemize}
  \item 它给出局部结果，即当前这一轮回应在中性度或边界维护上的即时反馈。
  \item 它也构成过程结果，因为多轮反馈会共同形成一次练习会话的轨迹。
\end{itemize}

与读心翻译相比，陪练路径更像对方法闭环的动态展开。系统不只判断一句话，也持续评估用户在压力语境中的回应方式是否保持克制、是否滑入解释或对抗。

\subsection{示例三：急救呼吸的状态恢复反馈}

急救呼吸路径的结果并不表现为文本建议，而表现为一个可观察的同步控制过程。页面当前可直接观察到 4-7-8 节奏、阶段文字同步与 1:59 倒计时。这里的结果属于状态型结果，即系统通过持续节奏引导帮助用户进入更稳定的呼吸与注意力节奏。

这一路径与前两条路径共同说明，Pebble 的结果面并不统一收敛为“文本答案”。在不同场景下，系统输出可以是结构化分析、回合反馈，也可以是恢复节奏本身。正因为存在这种差异，后文的验证也必须按路径分别处理，而不能用单一指标概括全系统。

\section{当前已观察结果}

若严格按照当前工作区中可直接观察到的事实来判断，Pebble 至少已经具备表~\ref{tab:pebble-observed-results} 所示的结果面。

\begin{table}[H]
  \centering
  \caption{Pebble 当前可直接观察到的结果类型及证据边界。}
  \label{tab:pebble-observed-results}
  \begin{tabularx}{\textwidth}{>{\raggedright\arraybackslash}p{2.7cm} >{\raggedright\arraybackslash}X >{\raggedright\arraybackslash}X}
    \toprule
    路径 & 当前可观察结果 & 证据边界 \\
    \midrule
    读心翻译 & 页面可展示表面语义、潜台词、情绪状态与三档回复建议，并支持复制与存入练习本 & 可确认结果结构已存在；不可仅据此宣称语义质量已充分验证 \\
    模拟陪练 & 页面可展示场景对话、AI 回复、情绪得分、即时反馈与注意点 & 可确认回合反馈链已建立；不可仅据此宣称训练效果已完成长期验证 \\
    急救呼吸 & 页面可展示 4-7-8 节奏、阶段文字、倒计时与暂停/重启交互 & 可确认恢复流程可运行；不可仅据此推断临床级稳定效果 \\
    练习沉淀 & 解码与陪练结果都可以进入练习材料对象 & 可确认沉淀入口存在；不可仅据此断言长期复盘价值已被量化 \\
    关系上下文 & 当前关系对象可穿透到翻译分析路径，形成上下文增强入口 & 可确认上下文注入链已存在；不可仅据此断言所有关系适配都足够精准 \\
    \bottomrule
  \end{tabularx}
\end{table}

上述结果的共同特点是，它们都属于“结构与行为已可观察”的层面。也就是说，当前可以确认系统已经形成若干稳定的输入输出形状、交互动作与状态流；但对这些结果是否已经达到高质量、强泛化或长期有效的程度，仍需更严格的验证证据支持。

\section{验证证据}

当前与 Pebble 相关的验证证据主要来自三类来源。

\begin{itemize}
  \item 页面与接口层的直接观察证据：当前工作区中的页面组件、状态管理与接口返回结构，能够直接证明结果形状和执行路径是否存在。
  \item 测试与前端断言证据：前端测试矩阵已覆盖翻译页、分析卡片、解码按钮、情绪状态条、回复建议与相关交互，具体覆盖范围与用例数量详见第 9 章"测试覆盖"小节。
  \item 任务书与测试计划证据：当前文档中已经给出了首页、翻译器、陪练场、急救呼吸与关系系统的验收项和测试步骤，为后续系统级验证提供了目标框架。已完成的执行清单与测试用例详见第 9 章"技术实现状态"与"测试覆盖"小节。
\end{itemize}

这三类证据的验证力度并不相同。

\begin{itemize}
  \item 直接观察证据最适合证明“是否存在某条链路或某种结果结构”。
  \item 前端测试最适合证明“在给定状态下，界面与局部交互是否按预期工作”。
  \item 验收方案最适合证明“系统原本想达到什么验收标准”，但它本身不等于已经执行完毕的结果。
\end{itemize}

因此，当前最稳妥的验证结论应停留在以下层次。

\begin{itemize}
  \item 可以确认 Pebble 的三条主路径都已经具备可观察的结果输出。
  \item 可以确认部分关键前端交互和显示结构已经进入测试覆盖。
  \item 可以确认系统已经存在较明确的验收框架与关系系统专项测试计划。
\end{itemize}

\section{尚未完成的验证边界}

本章必须同样明确哪些内容仍不能被夸大。

\begin{itemize}
  \item 当前并没有充分证据证明解码结果在广泛语境中都具有稳定高质量。
  \item 当前并没有充分证据证明模拟陪练已经形成可量化的长期训练收益。
  \item 当前并没有充分证据证明急救呼吸路径在临床或准临床场景中达到可外推的恢复效果。
  \item 当前也没有充分证据证明关系上下文增强已经在大量不同关系对象上得到系统性评估。
\end{itemize}

换言之，Pebble 当前最强的证据类型是“系统链路已存在并可运行”，而不是“长期效果已被严格量化”。这一边界需要在后续章节与最终交付中持续保持，否则报告很容易从工程事实滑向未经验证的效果宣称。

\section{本章小结}

本节读完后，读者应能区分 Pebble 当前已观察到的结果类型与尚不能验证的效果宣称，解释证据边界的理由，并列举至少两项仍需后续验证的系统能力。
